\chapter{Finger Pain}

\section*{Definition}
Finger pain is discomfort localized to one or more digits of the hand, resulting from trauma, arthritis, infection, or neuropathy. It is a common symptom that may impair fine motor function and grip strength.

\section*{Pathophysiology}
Pain-sensitive structures in the finger include the skin, subcutaneous tissue, periosteum, joint capsules, tendons, and digital nerves. Nociceptive pain arises from tissue injury (fracture, laceration, burn) activating A-delta and C fibers. Inflammatory arthritis involves synovial proliferation and cytokine release, while neuropathic pain results from direct nerve compression or injury (e.g., digital neuroma).

\section*{Causes}
\begin{itemize}
  \item \textbf{Traumatic:} Fractures (phalangeal, metacarpal), dislocations, lacerations, crush injuries, nail bed injuries, mallet finger, jersey finger, gamekeeper's thumb
  \item \textbf{Arthritic:} Osteoarthritis (Heberden nodes in DIP, Bouchard nodes in PIP), rheumatoid arthritis, psoriatic arthritis, gout, pseudogout
  \item \textbf{Infectious:} Felon (pulp space infection), paronychia, herpetic whitlow, flexor tenosynovitis
  \item \textbf{Neurologic:} Carpal tunnel syndrome, cervical radiculopathy, ulnar neuropathy
  \item \textbf{Vascular:} Raynaud phenomenon, digital ischemia, thromboangiitis obliterans
  \item \textbf{Other:} Ganglion cyst, trigger finger, Dupuytren contracture, foreign body
\end{itemize}

\section*{When to See a Doctor}
See your doctor if:
\begin{itemize}
  \item Severe pain after trauma with obvious deformity or open wound
  \item Signs of infection: erythema, warmth, purulent drainage, fever
  \item Inability to flex or extend the finger (tendon injury)
  \item Numbness or vascular compromise (pallor, cold, delayed capillary refill)
\end{itemize}

\section*{Related Symptoms}
\begin{itemize}
  \item Swelling and erythema of the affected digit
  \item Stiffness or reduced range of motion
  \item Numbness, tingling, or burning sensation
  \item Joint deformity (nodes, ulnar deviation, boutonniere, swan-neck)
  \item Crepitus with tendon movement in trigger finger
\end{itemize}
\textbf{Important:} Flexor tenosynovitis (Kanavel signs: fusiform swelling, flexed posture, tenderness along the tendon sheath, pain with passive extension) is a surgical emergency.

\section*{Self-Care / Home Management}
\begin{itemize}
  \item Rest, ice, compression, and elevation (RICE) for minor trauma
  \item Splinting of the affected finger in a functional position
  \item Over-the-counter NSAIDs or acetaminophen for pain
  \item Warm soaks for paronychia and minor infections (without fluctuance)
  \item Buddy taping for stable sprains
\end{itemize}

\section*{How Your Doctor Will Evaluate You}
\begin{itemize}
  \item Plain radiography (AP, lateral, and oblique views) for trauma or arthritis
  \item Ultrasound to evaluate tendon integrity, joint effusion, or foreign body
  \item MRI for occult fracture, osteomyelitis, or ligamentous injury
  \item Laboratory studies: CBC, ESR, CRP, uric acid, RF, anti-CCP if inflammatory arthritis suspected
  \item Wound culture and sensitivity for suspected infection
  \item Nerve conduction studies for suspected neuropathy
\end{itemize}

\section*{Treatment Options}
\begin{itemize}
  \item Reduction and immobilization for fractures and dislocations
  \item Antibiotics and incision and drainage for infections
  \item NSAIDs, corticosteroid injections, or disease-modifying antirheumatic drugs for arthritis
  \item Surgical repair of tendon lacerations or ruptures
  \item Physical and occupational therapy for range of motion and strengthening
\end{itemize}

\clearpage
